\documentclass[10pt,a4paper]{amsart}
\usepackage{amsmath,amssymb,amsthm,graphicx}
% Retained preamble for the Stacks Project chapter divisors.tex (Divisors).
% The chapter file pinned-original-source.tex carries only \input{preamble} on its first line, so no
% preamble text lives inside the chapter itself.  The Stacks Project uses one single shared preamble for
% every chapter; the snapshot stages it verbatim in sources/stacks-chow/preamble.tex.  The block below is
% copied mechanically, line for line, from that file's lines 14-18 (the Xy-pic setup) and lines 154-217
% (the native theorem declarations together with the mathematical macros).  Nothing was added, removed or
% redefined.  The wrapper class and the amsmath/amssymb/amsthm/graphicx base packages are supplied by the
% assembler, so the class/\externaldocument part of the shared preamble is not repeated here.
% For commutative diagrams we use Xy-pic
\usepackage[all]{xy}

% We use 2cell for 2-commutative diagrams.
\xyoption{2cell}
% Theorem environments.
%
\theoremstyle{plain}
\newtheorem{theorem}[subsection]{Theorem}
\newtheorem{proposition}[subsection]{Proposition}
\newtheorem{lemma}[subsection]{Lemma}

\theoremstyle{definition}
\newtheorem{definition}[subsection]{Definition}
\newtheorem{example}[subsection]{Example}
\newtheorem{exercise}[subsection]{Exercise}
\newtheorem{situation}[subsection]{Situation}

\theoremstyle{remark}
\newtheorem{remark}[subsection]{Remark}
\newtheorem{remarks}[subsection]{Remarks}

\numberwithin{equation}{subsection}

% Macros
%
\def\lim{\mathop{\mathrm{lim}}\nolimits}
\def\colim{\mathop{\mathrm{colim}}\nolimits}
\def\Spec{\mathop{\mathrm{Spec}}}
\def\Hom{\mathop{\mathrm{Hom}}\nolimits}
\def\Ext{\mathop{\mathrm{Ext}}\nolimits}
\def\SheafHom{\mathop{\mathcal{H}\!\mathit{om}}\nolimits}
\def\SheafExt{\mathop{\mathcal{E}\!\mathit{xt}}\nolimits}
\def\Sch{\mathit{Sch}}
\def\Mor{\mathop{\mathrm{Mor}}\nolimits}
\def\Ob{\mathop{\mathrm{Ob}}\nolimits}
\def\Sh{\mathop{\mathit{Sh}}\nolimits}
\def\NL{\mathop{N\!L}\nolimits}
\def\CH{\mathop{\mathrm{CH}}\nolimits}
\def\proetale{{pro\text{-}\acute{e}tale}}
\def\etale{{\acute{e}tale}}
\def\QCoh{\mathit{QCoh}}
\def\Ker{\mathop{\mathrm{Ker}}}
\def\Im{\mathop{\mathrm{Im}}}
\def\Coker{\mathop{\mathrm{Coker}}}
\def\Coim{\mathop{\mathrm{Coim}}}

% Boxtimes
%
\DeclareMathSymbol{\boxtimes}{\mathbin}{AMSa}{"02}

%
% Macros for moduli stacks/spaces
%
\def\QCohstack{\mathcal{QC}\!\mathit{oh}}
\def\Cohstack{\mathcal{C}\!\mathit{oh}}
\def\Spacesstack{\mathcal{S}\!\mathit{paces}}
\def\Quotfunctor{\mathrm{Quot}}
\def\Hilbfunctor{\mathrm{Hilb}}
\def\Curvesstack{\mathcal{C}\!\mathit{urves}}
\def\Polarizedstack{\mathcal{P}\!\mathit{olarized}}
\def\Complexesstack{\mathcal{C}\!\mathit{omplexes}}
% \Pic is the operator that assigns to X its picard group, usage \Pic(X)
% \Picardstack_{X/B} denotes the Picard stack of X over B
% \Picardfunctor_{X/B} denotes the Picard functor of X over B
\def\Pic{\mathop{\mathrm{Pic}}\nolimits}
\def\Picardstack{\mathcal{P}\!\mathit{ic}}
\def\Picardfunctor{\mathrm{Pic}}
\def\Deformationcategory{\mathcal{D}\!\mathit{ef}}

% --- appended 2026-10-12 (Chief): the project own definitions of the reference/slogan/history
% display environments, copied mechanically from the shared Stacks preamble lines 8-12
% (\usepackage{verbatim} plus the three \newenvironment lines, which the chapter preamble keeps
% verbatim from the project preamble). Nothing added, removed or redefined: these environments
% carry no mathematical content and the project itself typesets them as comments. Without these
% three lines a counted statement whose own text uses them fails with Environment slogan undefined.
\usepackage{verbatim}
\newenvironment{reference}{\comment}{\endcomment}
\newenvironment{slogan}{\comment}{\endcomment}
\newenvironment{history}{\comment}{\endcomment}
\usepackage{cleveref}
\crefname{originallemma}{original lemma}{original lemmas}
\Crefname{originallemma}{Original lemma}{Original lemmas}
\makeatletter
\newlabel{lemma-strict-transform-different-centers}{{080H}{source}{}{sourcepointer-1}{}}
\newlabel{lemma-strict-transform-different-centers@cref}{{[originallemma][0][]080H}{[1][0][]source}{}{}{}}
\makeatother
\title{The strict transform with respect to a composite blowing up is the iterated strict transform}
\author{The Stacks Project Authors}
\date{stacks-project:divisors.tex:local-snapshot-20261012; GNU FDL 1.2 or later}
\begin{document}
\maketitle
\noindent\textit{Editorial note.} All delivered mathematical text is the author's own, copied from the pinned Stacks Project chapter file pinned-original-source.tex; no statement or proof step was rewritten, reordered or removed.\par\medskip
\noindent\textit{Editorial note.} Following the multi-claim convention for this source, source\_url is the source-level URL https://stacks.math.columbia.edu/ and the source label lemma-strict-transform-composition-blowups is carried in claim\_key and here rather than in a per-claim tag URL.\par\medskip
\noindent\textit{Editorial note.} The preamble is the chapter's shared Stacks Project preamble, retained as sources/stacks-divisors/retained-preamble.tex: a line-for-line excerpt of the shared Stacks preamble staged with this snapshot. No macro was invented or redefined, and this card adds no preamble text of its own.\par\medskip
\noindent\textit{Editorial note.} Statement span 8969-8989 and proof span 8991-9034 come from the halfA pool card operations/pool-cards/stacks-divisors-halfA.json and were re-read in the pinned chapter file before this card was written.\par\medskip
\noindent\textit{Editorial note.} The proof contains a three-column xymatrix of base changes; it is copied verbatim and the retained chapter preamble loads Xy-pic.\par\medskip
\noindent\textit{Editorial note.} The statement assumes the composite blowing up presentation as a hypothesis rather than reproving it; the author's formulation is preserved.\par\medskip
\noindent\textit{Editorial note.} 1 cross-references occur inside the counted spans; they are listed in external\_references as navigation pointers only.\par\medskip
\noindent\textit{Allowed standard prerequisites (editorial).}\begin{itemize}
\item The definition of the strict transform of a quasi-coherent module and of a scheme with respect to a blowing up, by taking the image of the pullback of the module and the closure of the preimage of the complement of the centre. Used at proof lines 8991-8998 and 9021-9024. Supplies the objects being compared and the surjections that generate them.
\item The strict transform statement for two different centres, in the form that a strict transform acquires no nonzero sections supported on a prescribed centre given an effective Cartier divisor on it. Cited at proof lines 9030-9034. Shows that the iterated strict transform has no sections supported on the union of the two centres, which is the last step.
\item Right exactness of pullback, so that the two defining surjections compose into a single surjection from the pullback of the module. Used at proof lines 8999-9006. Produces the two surjections onto the iterated strict transform and onto the strict transform with respect to the composite whose kernels have to be compared.
\end{itemize}\medskip
\ifdefined\hypertarget\hypertarget{sourcepointer-1}{}\fi
\noindent\textit{Source reference.} Original lemma 080H: focused-seeds/AGP/sources/stacks-snapshot-20261011/tags.txt, source lines 10182--10182. Same-chapter reference used in the retained author text at source line(s) 9031. It is supplied only as a navigation pointer so that the retained text compiles; the counted item delivers the author's own statement and the author's complete proof, and no part of the referenced item is copied in or substituted for the counted argument.\par\medskip
\begin{lemma}
\label{lemma-strict-transform-composition-blowups}
Let $S$ be a scheme. Let $Z \subset S$ be a closed subscheme.
Let $b : S' \to S$ be the blowing up with center $Z$. Let $Z' \subset S'$ be
a closed subscheme. Let $S'' \to S'$ be the blowing up with center $Z'$.
Let $Y \subset S$ be a closed subscheme such that
$Y = Z \cup b(Z')$ set theoretically and the composition $S'' \to S$
is isomorphic to the blowing up of $S$ in $Y$.
In this situation, given any scheme $X$ over $S$ and
$\mathcal{F} \in \QCoh(\mathcal{O}_X)$ we have
\begin{enumerate}
\item the strict transform of $\mathcal{F}$ with respect to the blowing
up of $S$ in $Y$ is equal to the strict transform with respect to the
blowup $S'' \to S'$ in $Z'$ of the strict transform of $\mathcal{F}$
with respect to the blowup $S' \to S$ of $S$ in $Z$, and
\item the strict transform of $X$ with respect to the blowing
up of $S$ in $Y$ is equal to the strict transform with respect to the
blowup $S'' \to S'$ in $Z'$ of the strict transform of $X$
with respect to the blowup $S' \to S$ of $S$ in $Z$.
\end{enumerate}
\end{lemma}

\begin{proof}
Let $\mathcal{F}'$ be the strict transform of $\mathcal{F}$ with respect
to the blowup $S' \to S$ of $S$ in $Z$.
Let $\mathcal{F}''$ be the strict transform of $\mathcal{F}'$ with respect
to the blowup $S'' \to S'$ of $S'$ in $Z'$.
Let $\mathcal{G}$ be the strict transform of $\mathcal{F}$ with respect
to the blowup $S'' \to S$ of $S$ in $Y$.
We also label the morphisms
$$
\xymatrix{
X \times_S S'' \ar[r]_q \ar[d]^{f''} &
X \times_S S' \ar[r]_p \ar[d]^{f'} &
X \ar[d]^f \\
S'' \ar[r] & S' \ar[r] & S
}
$$
By definition there is a surjection $p^*\mathcal{F} \to \mathcal{F}'$
and a surjection $q^*\mathcal{F}' \to \mathcal{F}''$ which combine
by right exactness of $q^*$ to a surjection
$(p \circ q)^*\mathcal{F} \to \mathcal{F}''$. Also we have the surjection
$(p \circ q)^*\mathcal{F} \to \mathcal{G}$. Thus it suffices to prove
that these two surjections have the same kernel.

\medskip\noindent
The kernel of the surjection $p^*\mathcal{F} \to \mathcal{F}'$
is supported on $(f \circ p)^{-1}Z$, so this map is an isomorphism at
points in the complement. Hence the kernel of
$q^*p^*\mathcal{F} \to q^*\mathcal{F}'$
is supported on $(f \circ p \circ q)^{-1}Z$. The kernel of
$q^*\mathcal{F}' \to \mathcal{F}''$ is supported on $(f' \circ q)^{-1}Z'$.
Combined we see that the kernel of
$(p \circ q)^*\mathcal{F} \to \mathcal{F}''$ is supported on
$(f \circ p \circ q)^{-1}Z \cup (f' \circ q)^{-1}Z' =
(f \circ p \circ q)^{-1}Y$.
By construction of $\mathcal{G}$ we see that we obtain a factorization
$(p \circ q)^*\mathcal{F} \to \mathcal{F}'' \to \mathcal{G}$.
To finish the proof it suffices to show that $\mathcal{F}''$ has no
nonzero (local) sections supported on
$(f \circ p \circ q)^{-1}(Y) =
(f \circ p \circ q)^{-1}Z \cup (f' \circ q)^{-1}Z'$.
This follows from Lemma \ref{lemma-strict-transform-different-centers}
applied to $\mathcal{F}'$ on $X \times_S S'$ over $S'$, the closed
subscheme $Z'$ and the effective Cartier divisor $b^{-1}Z$.
\end{proof}

\end{document}
